\documentclass[10pt]{amsart}
\usepackage[a4paper,margin=23mm]{geometry}
\usepackage{amsmath,amssymb,amsthm,mathtools,hyperref,graphicx,xspace,letltxmacro,xpatch}
\usepackage[english]{babel}
\newtheorem{theo}{Theorem}[section]
\newtheorem{lemm}[theo]{Lemma}\newtheorem{prop}[theo]{Proposition}\newtheorem{coro}[theo]{Corollary}\newtheorem{conj}[theo]{Conjecture}
\theoremstyle{definition}\newtheorem{defi}[theo]{Definition}\newtheorem{rema}[theo]{Remark}\newtheorem{exam}[theo]{Example}\newtheorem{exem}[theo]{Example}\newtheorem{nota}[theo]{Notation}
\newcommand{\ie}{i.e.}\newcommand{\eg}{e.g.}\newcommand{\resp}{resp.}\newcommand{\skpt}{\par\smallskip\noindent}
\emergencystretch=3em\pagestyle{plain}

%%\usepackage{mathtools}
%
\usepackage{bm}






\newcommand{\Fca}{\mathcal{F}}
\newcommand{\Gca}{\mathcal{G}}
\newcommand{\Hca}{\mathcal{H}}
\newcommand{\Rca}{\mathcal{R}}
\newcommand{\clH}{\mathcal{H}}
\newcommand{\clB}{\mathcal{B}}
\newcommand{\f}{\mathcal{F}}

\newcommand{\bfun}{\mathbf{1}}

\newcommand{\rmd}{\mathrm{d}}
\newcommand{\seq}{\mathrm{seq}}

\newcommand\R{\mathbb{R}}
\newcommand\N{\mathbb{N}}
\newcommand\PP{\mathbb{P}}
\newcommand\Z{\mathbb{Z}}
\newcommand{\bbE}{\mathbb{E}}
\newcommand{\bbL}{\mathbb{L}}

\newcommand{\E}[1]{\mathbb E[#1]}
\newcommand\ens[1]{\{ #1\}}
\newcommand\abs[1]{|#1|}
\newcommand{\pr}[1]{(#1)}
\newcommand{\norm}[1]{\lVert #1\rVert}

\newcommand\eps{\varepsilon}

\newcommand{\gr}[1]{\bm{#1}}
\newcommand{\gri}{\bm{i}}
\newcommand{\grj}{\bm{j}}
\newcommand{\grn}{\bm{n}}
\newcommand{\grm}{\bm{m}}
\newcommand{\gru}{\bm{u}}
\newcommand{\grk}{\bm{k}}
\newcommand{\grt}{\bm{t}}
\newcommand{\imd}{\preccurlyeq}
\newcommand{\smd}{\succcurlyeq }

\makeatletter
\def\editors#1{%
\def\editor@name{#1}
\if@francais
\def\editor@string{Recommandé par les éditeurs \editor@name.}
\else
\def\editor@string{Recommended by Editors \editor@name.}
\fi}
\makeatother

%%%%%%%%%%%%%%%%%%%%%%%%%%%%%%%%%%%%%%%%%%%%%%%%%%
\graphicspath{{./figures/}}

\newcommand*{\mk}{\mkern -1mu}
\newcommand*{\Mk}{\mkern -2mu}
\newcommand*{\mK}{\mkern 1mu}
\newcommand*{\MK}{\mkern 2mu}

%\hypersetup{urlcolor=purple, linkcolor=blue, citecolor=red}

\newcommand*{\romanenumi}{\renewcommand*{\theenumi}{\roman{enumi}}}
\newcommand*{\Romanenumi}{\renewcommand*{\theenumi}{\Roman{enumi}}}
\newcommand*{\alphenumi}{\renewcommand*{\theenumi}{\alph{enumi}}}
\newcommand*{\Alphenumi}{\renewcommand*{\theenumi}{\Alph{enumi}}}
\newcommand*{\stepenumi}{\renewcommand*{\theenumi}{Step~\arabic{enumi}}}
\let\oldtilde\tilde
\renewcommand*{\tilde}[1]{\mathchoice{\widetilde{#1}}{\widetilde{#1}}{\oldtilde{#1}}{\oldtilde{#1}}}
\let\oldforall\forall
\renewcommand*{\forall}{\mathrel{\oldforall}}
%%%%%%%%%%%%%%%%%%%%%%%%%%%%%%%%%%%%%%%%%%%%%	

\title{Exponential rectangular-sum tail bounds for orthomartingale fields with shift-invariant block laws}
\author{Giraudo, Davide}
\date{doi:10.5802/ahl.172; CC BY 4.0}
\begin{document}
\maketitle
\noindent\textit{Editorial note.} Only selected main theorem counted; entire original mathematical body, exact definitions and proof context retained. Native bibliography keys preserved. Standard prerequisites remain author-cited and are not recursively reproved.\par\medskip
\noindent\textit{Allowed standard prerequisites (editorial).}\begin{itemize}
\item Standard functional analysis, measure theory and the author-cited classical theorems. Standard prior results stay cited; complete source-local proofs and definitions retained.
\end{itemize}\medskip

\section{Statement of the exponential inequality}

\subsection{Goal and motivations} Understanding the asymptotic behavior of sums of random variables is an important topic in probability. A useful tool for establishing limit theorems is to control the probability that the partial sums or the maximum of the absolute values of the partial sums exceeds a fixed number. A direct application can give convergence rates in the law of large numbers, and such an inequality can be used to check tightness criteria in some functional spaces.

In this paper, we will be concerned by establishing and applying an exponential inequality for the so-called orthomartingale difference random fields introduced in~\cite{MR0254912}. We make an assumption of invariance in law of the partial sums, which is weaker than stationarity. We control the tails of maximum of partial sums on rectangles of $\N^d$ by two quantities: an exponential term and another one involving the common distribution function of the increments. We will then formulate two applications of the result. One will deal with the convergence rates in the strong law of large numbers. The second one will be about the invariance principle in H\"older spaces, which will be restricted to the case of stationary processes for technical reasons.

Orthomartingales are well-suited for the summation over rectangles of $\Z^d$ because we can use a one dimensional martingale property when we sum according to each fixed coordinate. Moreover, the approximation of stationary random fields by this class of martingales can be done, like in~\cite{MR3504508,MR3869881,MR3798239,MR4166203}. However, for rates in the law of large numbers and the functional central limit theorem in H\"older spaces, only a few results are available in the literature. Indeed, rates on the law of large number for orthomartingales with polynomial moment have been given in~\cite{MR4046858,MR2794415,MR3451971}, but it seems that the question of exponential moments was only addressed for martingale difference sequences. For the functional central limit theorem in H\"older spaces for random fields, the i.i.d.\ case was addressed for moduli of regularity of the form $t^\alpha$, $0<\alpha<1/2$ in the case of random fields, and for more general moduli in the case of sequences. It turns out that the inequality we will present in this paper is a well-suited tool for dealing with these problems.

A key ingredient for the aforementioned limit theorems is an exponential inequality for orthomartingale random fields, that is, an inequality putting into play an exponential term and the tail of the common distribution of the increments. At first glance, it seems that multi-indexed martingales could be treated like sequences, up to some technical and notational obstacles. However, it turns out that the standard tools for proving deviation inequalities for martingales, like martingale transforms or uses of exponential supermartingale, do not extend easily. Nevertheless, it is possible to apply induction arguments when the random field satisfies good properties.

The paper is organized as follows: in Subsection~\ref{subsec:expo_ineg}, we state the definition of orthomartingale difference random fields and an exponential inequality for such random fields, with an assumption of invariance of the law of the sum on rectangles by translation. Subsections~\ref{subsec:LLN} (respectively~\ref{subsc:HWIP}) provide an application to the rates in the strong law of large numbers (respectively, to the invariance principle in H\"older spaces). Section~\ref{sec:proofs} is devoted to the proof of the previously stated results.

\subsection{Exponential inequality for orthomartingales}\label{subsec:expo_ineg}

We start by defining the concept of orthomartingale. Before that, we need to introduce the notion of filtration with respect to the coordinatewise order $\imd$: we say that for $\gri,\grj\in\Z^d$, $\gri\imd\grj$ if $i_q\leq j_q$ for each $q\in \ens{1,\,\dots,\,d}=:[d]$.

\begin{defi}
A filtration indexed by $\Z^d$ on a probability space $\pr{\Omega,\f,\PP}$ is a collection of $\sigma$-algebras $\pr{\f_{\gri}}_{\gri\in \Z}$ such that for each $\gri\imd\grj$, the inclusion $\f_{\gri}\subset
\f_{\grj}$ takes place.
\end{defi}

\begin{defi}
We say that the filtration $\pr{\f_{\gri}}_{\gri\,\in\,\Z}$ is commuting if for all integrable random variable $Y$ and all $\gri$, $\grj\in\Z$, the equality
\begin{equation}
\bbE\left[\left[Y\middle| \f_{\gri}\right]\middle|\f_{\grj}\right] = \bbE\left[Y\middle| \f_{\min\left\{\gri,\grj\right\}}\right],
\end{equation}
takes place, where the minimum is taken coordinatewise.
\end{defi}

\begin{exam}
Let $\pr{\eps_{\grj}}_{\grj\,\in\,\Z^d}$ be an i.i.d.\ random field and let $\f_{\grj}:=\sigma\pr{\eps_{\gri},\gri\imd\grj}$. Then the filtration $\pr{\f_{\gri}}_{\gri\,\in\,\Z}$ is commuting.
\end{exam}

\begin{exam}
Let $\pr{\eps_k^{\pr{q}}}_{k\,\in\,\Z}$, $q\in [d]:=\ens{1,\,\dots,\,d}$ be independent copies of an i.i.d.\ sequence $\pr{\eps_k}_{k\,\in\,\Z}$ and let $\f_{\grj}:=\sigma\pr{\eps_{k_q}^{\pr{q}},k_q\leq j_q,1\leq q\leq d}$. Then the filtration $\pr{\f_{\gri}}_{\gri\,\in\,\Z}$ is commuting, and connected to decoupled $U$-statistics.
\end{exam}


\begin{defi}
We say that the random field $\pr{X_{\gri}}_{\gri\,\in\,\Z^d}$ is an orthomartingale difference random field with respect to the commuting filtration $\pr{\f_{\gri}}_{\gri\,\in\,\Z^d}$ if for all $\gri\in\Z^d$, $X_{\gri}$ is integrable, $\f_{\gri}$-measurable and
\begin{equation}
\left[X_{\gri}\middle| \f_{\gri-\gr{e}_q}\right]=0,
\end{equation}
where for $q\in [d]$, $\gr{e}_q$ is the $q^{\rm th}$ vector of the canonical basis of $\R^d$.
\end{defi}



\begin{theo}\label{thm:deviation_inequality_orthomartingale}
Let $\pr{X_{\gri}}_{\gri\,\in\,\Z^d}$ be an orthomartingale differences random field such that for all $\grn\in\N^d$ and all $\grk\in\Z^d$,
\begin{equation}\label{eq:invariance_loi_sommes}
S_{\grn}:=\sum_{\gr{1}\,\imd\,\gri\,\imd\,\grn} X_{\gri}\text{ and } \sum_{\gr{1}\,\imd\,\gri\,\imd\,\grn} X_{\gri+\grk}\text{ have the same distribution}.
\end{equation}
Then the following inequality holds for all $x,y>0$:
\begin{multline}\label{eq:exp_deviation_inequality_orthomartingale}
\PP\left\{\max_{\gr{1}\,\imd\,\gri\,\imd\,\gr{n}} \left|S_{\gri} \right| >x\left|\grn\right|^{1/2}\right\}\leq A_d\exp\left(-\left(\frac{x}{y}\right)^{2/d}\right)
\\+B_d\int_1^{+\infty}\PP\left\{\left|X_{\gr{1}}\right|>yuC_d\right\} u\left(\log\pr{1+u}\right)^{p_d}du,
\end{multline}
where $A_d$, $B_d$ and $C_d$ depend only on $d$, $p_d=2d$ and $\abs{\grn}=\prod_{q=1}^dn_q$.
\end{theo}


\begin{rema}
Assume that $\pr{X_{\gri}}_{\gri\,\in\,\Z^d}$ is bounded, that is, there exists a constant $K$ such that $\abs{X_{\gr{i}}}\leq K$ almost surely for all $\gri\in\Z^d$. If $x>3^{d/2}K/C_d$ then we can choose $y=K/C_d$ and inequality~\eqref{eq:exp_deviation_inequality_orthomartingale} simplifies as
\begin{equation}
\PP\left\{\max_{\gr{1}\,\imd\,\gri\,\imd\,\gr{n}} \left|S_{\gri} \right| >x\left|\grn\right|^{1/2}\right\}\leq A_d\exp\left(-\left(\frac{C_dx}{K}\right)^{2/d}\right).
\end{equation}
\end{rema}


\begin{rema}
The exponent $2/d$ in the exponential term of~\eqref{eq:exp_deviation_inequality_orthomartingale} is not improvable, even in the bounded case. To see this, consider $d$ i.i.d. sequences $\pr{\eps_i^{\pr{q}}}_{i\,\in\,\Z}$ which are independent of each other, in the sense that the collection $\ens{\eps_{i_q}^{\pr{q}}, i_q\,\in\,\Z}$ is independent. Assume that $\eps_0^{\pr{q}}$ takes the values $1$ and $-1$ with probability $1/2$ for all $q\in [d]$. Let $X_{\gri}=\prod_{q=1}^d\eps_{i_q}^{\pr{q}}$. Then
\begin{multline}
\PP\left\{\max_{\gr{1}\,\imd\,\gri\,\imd\,\gr{n} } \left|S_{\gri} \right| >x\left|\grn\right|^{1/2}\right\}\geq
\PP\left\{ \left|S_{\grn} \right| >x\left|\grn\right|^{1/2}\right\}\\
=
\PP\left\{ \left|\prod_{q=1}^d
\frac 1{\sqrt{n_q}}\sum_{i_q=1}^{n_q}\eps_{i_q}^{\pr{q}} \right| >x \right\}.
\end{multline}
The vector $\pr{\frac 1{\sqrt{n_q}}\sum_{i_q=1}^{n_q}\eps_{i_q}^{\pr{q}}}_{q=1}^d$ converges in distribution to $\pr{N_q}_{q=1}^d$, where $\pr{N_q}_{q=1}^d$ is independent and each $N_q$ has a standard normal distribution. Therefore, if $f\pr{x}$ is a function such that for each bounded by $1$ orthomartingale difference random field satisfying~\eqref{eq:invariance_loi_sommes} and each $x>0$
\begin{equation}
\PP\left\{\max_{\gr{1}\,\imd\,\gri\,\imd\,\gr{n}} \left|S_{\gri} \right| >x\left|\grn\right|^{1/2}\right\}\leq f\pr{x},
\end{equation}
then letting $Y:= \prod_{q=1}^d\abs{N_q}$, the following inequality should hold for all $x$: $\PP\ens{Y>x}\linebreak\leq f\pr{x}$. Since the $\bbL^p$-norm of $Y$ behaves like $p^{d/2}$, the function $f$ cannot decay quicker than $\exp\pr{-Kx^{\gamma}}$ for some $K>0$ and $\gamma>2/d$.
\end{rema}

\begin{rema}
The $\log$ factor in the right hand side of~\eqref{eq:exp_deviation_inequality_orthomartingale} appears naturally as iterations of weak-type estimates of the form $x\PP\ens{X>x}\leq \E{X\bfun\ens{Y>x}}$ for some random variable $Y$, giving a control of the tail of $X$ in terms of that of $Y$. Consequently, the $\log$ factor does not seem to be avoidable with this method of proof. We do not know whether this factor can be removed.
\end{rema}


\section{Application to limit theorems}

\subsection{Convergence rates in the law of large numbers}\label{subsec:LLN}
A centered sequence $\pr{X_i}_{i\,\geq\,1}$ satisfies the law of large numbers if the sequence $\pr{n^{-1}\sum_{i=1}^nX_i}_{i\,\geq\,1}$ converges almost surely to zero. This is for example the case of a strictly stationary ergodic sequence where $X_1$ is integrable and centered. Then arises the question of evaluating the speed of convergence and finding bounds for the large deviation probabilities, namely,
\begin{equation}
\PP\left\{\frac 1n\left| \sum_{i=1}^nX_i\right| >x\right\}.
\end{equation}
This question has been treated in the independent case in~\cite{MR30714,MR19852,MR79851} under conditions on the $\bbL^p$-norm of $X_i$. The case of martingale differences has also been addressed: under boundedness of moments of order $p$ and a Cramér-type condition $\sup_{i\,\geq\,1}\E{\exp\pr{\abs{X_i}}}<+\infty$ in~\cite{MR1856684}, a conditional Cramér condition in~\cite{MR2568267} and under finite exponential moments in~\cite{MR3005732}.

For random fields, we can consider large deviation probabilities defined by
\begin{equation}
\PP\left\{\frac 1{\left|\gr{N}\right|}\left| \sum_{\gr{1}\,\imd
\,\grn\,\imd\,\gr{N}} X_{\gri}\right| >x\right\}, \gr{N} \smd\gr{1}, x>0.
\end{equation}
Result for orthomartingales with polynomial moment have been given in~\cite{MR4046858,MR2794415,MR3451971}.

\begin{theo}\label{thm:large_deviation}
Let $\pr{X_{\gri}}_{\gri\,\in\,\Z^d}$ be an orthomartingale difference random field satisfying~\eqref{eq:invariance_loi_sommes}. Suppose that for some $\gamma>0$,
\begin{equation}
\sup_{s\,>\,0}\exp\left(s^\gamma\right)\PP\left\{\left|X_{\gr{1}}\right|>s\right\}\leq 2.
\end{equation}
Then for each positive $x$, the following inequality takes place:
\begin{equation}
\PP\left\{\frac 1{\left|\gr{N}\right|}\left| \sum_{\gr{1}\,\imd\,
\gr{n}\,\imd\,\gr{N}} X_{\gri}\right| >x\right\}\leq C_{1,d,\gamma}\exp\left(-C_{2,d,\gamma}
\left|\gr{N}\right|^{\frac{\gamma}{2+d\gamma}}x^{\frac{2\gamma}{2+d\gamma}} \right),
\end{equation}
where $C_{1,d,\gamma}$ and $C_{2,d,\gamma}$ depend only on $d$ and $\gamma$.
\end{theo}



\subsection{Application to H\"olderian invariance principle}\label{subsc:HWIP}
Given a sequence of random variables $\pr{X_i}_{i\,\geq\,1}$, a way to understand the asymptotic behavior of the partial sums given by $S_k=\sum_{i=1}^kX_i$ is to define a sequence of random functions $\pr{W_n}_{n\,\geq\,1}$ on the unit interval $[0,1]$ in the following way: $W_n\pr{k/n} =n^{-1/2}S_k$ for $k\in\ens{1,\,\dots,\,n}$, $W_n\pr{0}=0$ and on $\pr{\pr{k-1}/n,k/n}$, $W_n$ is linearly interpolated. When the sequence $\pr{X_i}_{i\,\geq\,1}$ is i.i.d., centered and has unit variance, the sequence $\pr{W_n}_{n\,\geq\,1}$ converges in law in $C[0,1]$ to a standard Brownian motion. The result has been extended to strictly stationary martingale difference sequences in~\cite{MR0126871,MR0151997}. Then numerous papers treated the case of weakly dependent strictly stationary sequence, see~\cite{MR2206313} and the references therein for an overview.

There are also two other possibilities of extension of such result. The first one is to consider other functional spaces, in order to establish the convergence of $\pr{F\pr{W_n}}_{n\,\geq\,1}$ for a larger class of functionals than the continuous functionals $F\colon C[0,1]\to\R$. In other words, we view $W_n$ as an element of an element of a H\"older space $\clH_\rho$ with modulus of regularity $\rho$ and investigate the convergence of $\pr{W_n}_{n\,\geq\,1}$ in this function space. Applications to epidemic changes can be given, for example in~\cite{MR2265057}. The question of the invariance principle in H\"older spaces has been treated for i.i.d.\ sequences for modulus of regularity of the form $t^\alpha$ in~\cite{MR2000642} and also for more general ones, of the form $t^{1/2}\log
\pr{ct}^\beta$ in~\cite{MR2054586}. Some results are also available for stationary weakly dependent sequences, for example mixing sequences~\cite{MR3615086,MR1759810} or by the use of a martingale approximation~\cite{MR3426520,MR3770868}.

A second development of these invariance principles is the consideration of partial sum processes built on random fields. Given a random field $\pr{X_{\gri}}_{\gri\,\in\,\Z^d}$, one can define the following random function on $[0,1]^d$:
\begin{equation}
W_{\grn}\left(\grt\right)=\frac 1{\sqrt{\left|\grn\right|}}
\sum_{\gri\,\in\,\Z^d}
\lambda\left(R_{\gri}\cap \prod_{q=1}^d \left[0,n_qt_q\right]\right) X_{\gri}, \grn\smd\gr{1}, \grt \in [0,1]^d,
\end{equation}
where $\lambda$ denotes the Lebesgue measure on $\R^d$ and $R_{\gri}=\prod_{q=1}^d\pr{ i_q-1,i_q}$. Notice that the map $\grt\mapsto W_{\grn}\pr{\grt}$ is $1$-Lipschitz continuous and that $W_{\grn}\pr{\pr{\frac{k_q}{n_q}}_{q=1}^d}=S_{k_1,\,\dots,\,k_d}/\sqrt{\abs{\grn}}$ hence $W_{\grn}$ takes into account the values of all the partial sums $S_{\grk}$ for $\gr{1}\imd\gr{k}\imd\grn$.

Convergence of $W_{\grn}$ in the space of continuous functions has been investigated in~\cite{MR0246359} for an i.i.d.\ random field, \cite{MR1875665} for martingales for the lexicographic order and in~\cite{MR3427925,MR3913270} for orthomartingales.

In this section, we will study the convergence of $W_{\grn}$ in some H\"older spaces for orthomartingales random fields. First, an observation is that for $\grt,\gr{t'}\in [0,1]^d$,
\begin{multline}
\left| \lambda\left(R_{\gri}\cap \prod_{q=1}^d [0,n_qt_q]\right)- \lambda\left(R_{\gri}\cap \prod_{q=1}^d [0,n_qt'_q]\right)\right|\\
\leq \max_{1\,\leq\,q\,\leq\,d}n_q
\left|t_q-t'q\right|\leq \left|\grn\right|\left\|\grt-\gr{t'}\right\|_\infty
\end{multline}
hence
\begin{equation}
\left|W_{\grn}\left(\grt\right)-W_{\grn}\left(\gr{t'}\right)\right|
\leq \sqrt{\left|\grn\right|} \left\|\grt-\gr{t'}\right\|_\infty
\sum_{\gr{1}\,\imd\,\gri\,\imd\,\grn}\left|X_i\right|
\end{equation}
hence for almost every $\omega$, the map $\grt\mapsto W_{\grn}\pr{\grt}\pr{\omega}$ is Lipschitz-continuous. As pointed out in the i.i.d.\ case in~\cite{MR0246359}, the finite dimensional distributions of $W_{\grn}$ converge to those of a standard Brownian sheet, that is, a centered Gaussian process $\pr{W\pr{\grt}}_{\grt\,
\in\,[0,1]^d}$ whose covariance function is given by
\begin{equation}
\operatorname{Cov}\big(W\pr{\grt},W\pr{\gr{t'}}\big) =\prod_{q=1}^d\min\left\{t_q,t'_q\right\},
\end{equation}
provided that $X_{\gr{1}}$ is centered and has unit variance. Given an increasing continuous function $\rho\colon [0,1]\to\R$ such that $\rho\pr{0}=0$, let $\Hca_\rho$ be the space of all the functions $x$ from $[0,1]^d$ to $\R$ such that $\sup_{\grt,\gr{t'}\,\in\,[0,1]^d,
\grt\neq\gr{t'}}\abs{x\pr{\grt}-x\pr{\gr{t'}}}/
\rho\pr{\norm{\grt-\gr{t'}}_\infty }$ is finite.

The Brownian sheet has trajectories in $\Hca_\rho$ with $\rho\pr{h} =h^\alpha$ for all $\alpha \in\pr{0,1/2}$ but not for $\alpha\geq 1/2$. Therefore, it is not possible to expect to show the convergence of $\pr{W_{\grn}}_{\grn\,\smd\,\gr{1}}$ in all the possible H\"older spaces and some restriction have to be made. In order to treat a class of modulus of regularity larger than the power function, we need to introduce the slowly varying functions. We say that a function $L\colon \pr{0,+\infty}\to \pr{0,+\infty}$ is slowly varying if for all $c>0$, the quantity $L\pr{ct}/L\pr{t}$ goes to $1$ as $t$ goes to infinity. For example, functions which behave asymptotically as a power of the logarithm are slowly varying. We now define the moduli of regularity and the associated H\"older spaces of interest from the point of view of the convergence of the partial sum process $W_{\grn}$.

\begin{defi}
Let $d\geq 1$ be an integer. We say that $\rho$ belongs to the class $\Rca_{1/2,d}$ if there exists a slowly varying function $L$ such that $L\pr{t}\to \infty$ as $t$ goes to infinity and a constant $c$ such that
\begin{equation}\label{eq:def_de_rho}
\rho\pr{h}=h^{1/2}\left(\ln\left(\frac{c}{h}\right)\right)^{d/2}L\left(\frac 1h\right), h
\in [0,1],
\end{equation}
and $\rho $ is increasing on $[0,1]$.
\end{defi}

It seems that the exponent $d/2$ is the best we can get in view of proving tightness via the deviation inequality we established. It may not be optimal in some cases, for example, if $\pr{X_{\gri}}_{\gri\,\in\,\Z^d}$ is centered, we can get a similar inequality as~\eqref{eq:exp_deviation_inequality_orthomartingale}, but with the exponent $2$ instead of $2/d$ in the right hand side. As a consequence, tightness in $\Hca_{\rho}$ can be established for $\rho\in\Rca_{1/2,1}$, a larger class than $\rho\in\Rca_{1/2,d}$


We now give a sufficient condition for tightness of the partial sum process associated to a strictly stationary random field, that is, a random field $\pr{X_{\gri}}_{\gri\,\in\,\Z^d}$ such that for each integer $N$ and each $\gr{i_1},\,\dots,\,\gr{i_N},\gr{j}\in\Z^d$, the vectors $\pr{X_{\gr{i_k}+\grj}}_{k=1}^N$ and $\pr{X_{\gr{i_k}}}_{k=1}^N$ have the same distribution. The criterion puts into play tails of the maximum of partial sums of the rectangles.

\begin{prop}\label{prop:tightness_partial_sums}
Let $d\geq 1$ and $\rho$ be an element of $\Rca_{1/2,d}$. If $\pr{X_{\gri}}_{\gri\,\in\,\Z^d}$ is a strictly stationary random field such that for each $q\in \ens{1,\,\dots,\,d}$ and each positive~$\varepsilon$,
\begin{equation}\label{eq:tightness_random_fields_sufficient_cond}
\lim_{J\,\to\,\infty}\limsup_{\min\,\grm\,\to\,\infty}
\sum_{j=J}^{m_q}2^j \PP\left\{\max_{\gr{1}\,\imd\,\gr{k}\,\imd\,2^{\grm -m_q\gr{e}_q } }
\left|\sum_{\gr{1}\,\imd\,\gri\,\imd\,\grk} X_{\gri} \right|> \varepsilon\rho\left(2^{-j}\right)
\prod_{u=1}^d2^{m_u/2}\right\}=0,
\end{equation}
then the net $(W_{\grn})_{\grn\,\smd\, \gr{1}}$ is asymptotically tight in $\clH_{\rho}([0,1]^d)$ as $\min\grn=
\min_{1\,\leq\,q\,\leq\,d}n_q\linebreak\to\infty$.
\end{prop}

It turns out that it is more convenient to consider maximum indexed by dyadic elements of $\Z^d$. In order to avoid confusion, we will use notations of the form $2^{m_k}$ instead of $2^{n_k}$.


A similar tightness criterion was used in~\cite{MR3615086} for sequences, giving optimal results for mixing sequences and allowing to recover the optimal result for i.i.d.\ sequences. For random fields, the inequality~\eqref{eq:exp_deviation_inequality_orthomartingale} we obtained and that in Theorem~1.13 of~\cite{MR4046858} are appropriate tools to check~\eqref{eq:tightness_random_fields_sufficient_cond}.


No other assumption than stationarity is done but of course, some dependence will be required for this condition to be satisfied since one need a good control of the tails of the partial sums on rectangles normalized by the square root of the number of elements in the rectangle.

Now, we state a result for the weak convergence of $\pr{W_{\grn}}_{\grn\,\smd\,\gr{1}}$ in the space $H_\rho$, where $\rho\in\Rca_{1/2,d}$.

\begin{theo}\label{thm:WIP_Holder}
Let $d\geq 1$ and let $\pr{X_{\gri}}_{\gri\,\in\,\Z^d}$ be a strictly stationary orthomartingale difference random field. Let $\rho$ be an element of $\Rca_{1/2,d}$ be given by~\eqref{eq:def_de_rho}, where $L$ is slowly varying.

Suppose that
\begin{equation}\label{eq:cond_suff_WIP}
\forall A>0, \quad\sum_{j\,\geq\,1}2^j\PP\left\{\left|X_{\gr{1}}\right|>L\left(2^j\right)A\right\}<\infty.
\end{equation}
Then $\pr{W_{\grn}}_{\grn\,\smd\,\gr{1}}$ is asymptotically tight in $\Hca_\rho\pr{[0,1]^d}$ as $\min\grn=
\min_{1\,\leq\,q\,\leq\,d}n_q\to\infty$.

Assume moreover that any $A\in\clB\pr{\R^{\Z^d}}$ such that 
\[
\left\{\left(X_{\gri}\right)_{\gri\,\in\,\Z^d}\in A\right\}=\left\{\left(X_{\gri+\gr{e}_d}\right)_{\gri\,\in\,\Z^d}\in A\right\}
\]
satisfies $\PP\ens{\pr{X_{\gri}}_{\gri\,\in\,\Z^d}\in A}\in\ens{0,1}$. Then $\pr{W_{\grn}}_{\grn\,\smd\,\gr{1}}$ converges to $\norm{X_{\gr{0}}}_2 W$ in $\Hca_\rho\linebreak\pr{[0,1]^d}$ as $\min\grn=
\min_{1\,\leq\,q\,\leq\,d}n_q\to\infty$, where $W$ is a standard Brownian sheet.
\end{theo}


\section{Proofs}\label{sec:proofs}


\subsection{Proof of Theorem~\ref{thm:deviation_inequality_orthomartingale}}

The proof will be done by induction on the dimension $d$. For $d=1$, we need the following deviation inequality for one dimensional martingales. This is a combination of Theorem~2.1 in~\cite{MR3311214} and Theorem~6 in~\cite{MR606989}.

\begin{prop}[Proposition~2.1 in
\cite{MR4466386}]\label{prop:inegalite_deviation_martingales_dominee}
Let $\pr{D_i}_{i\,\geq\,1}$ be a martingale difference sequence with respect to the filtration $\pr{\Fca_i}_{i\,\geq\,0}$. Suppose that $\E{D_i ^2}$ is finite for all $i\geq 1$. Suppose that there exists a nonnegative random variable $Y$ such that for all $1\leq i\leq n$, $\E{\varphi\pr{D_i^2}}\leq\E{\varphi\pr{Y^2}}$ for all convex increasing function $\varphi\colon [0,\infty)\to[0,\infty)$. Then for all $x,y>0$ and each $n\geq 1$, the following inequality holds:
\begin{multline}\label{eq:inegalite_max_martingales_moments_ordre_p_domination_sto}
\PP\left\{\max_{1\,\leq\,k\,\leq\,n}\left|\sum_{i=1}^kD_i\right|>x n^{1/2} \right\}
\\
\leq 2\exp\left(-\frac 12\left(\frac{x}{y}\right)^{2} \right)
+ 4\int_{1 }^{+\infty}w\PP\left\{ Y>yw/2\right\}\rmd w.
\end{multline}
\end{prop}

\begin{proof}
Assume that Theorem~\ref{thm:deviation_inequality_orthomartingale} holds in dimension $d-1$ with $d\geq 2$. We have to prove~\eqref{eq:exp_deviation_inequality_orthomartingale} for all $d$-dimensional orthomartingale difference fields satisfying~\eqref{eq:invariance_loi_sommes}. We first get rid of the maximum over the coordinate $d$, apply the inequality of the $d-1$ dimension to the obtained orthomartingale and we are then reduced to control the tails of a one dimensional martingale. More concretely, the induction step is done as follows.
\begin{enumerate}\stepenumi
\item\label{step1} let $M:=\max_{\gr{1}\,\imd\,\gri\,\imd\,\gr{n} } \abs{S_{\gri} }$ and
\begin{equation}
M':=\max_{1\,\leq\,i_q\,\leq\,n_q,\,1\,\leq\,q\leq\,d-1}
\left|S_{i_1,\,\dots,\,i_{d-1},\,n_d}\right|.
\end{equation}
Then we will show that
\begin{equation}
\PP\left\{M>x\right\}\leq \int_1^\infty\PP\left\{M'>xu/2\right\}du.
\end{equation}
\item\label{step2} the tails of $M'$ are controlled by applying the result for $\pr{d-1}$-dimensional random fields.
\item\label{step3} it remains to control the tails of $\sum_{i_d=1}^{n_d}X_{1,\,\dots,\,1,i_d}$, which can be done by using the one dimensional result.
\end{enumerate}


The proof will be quite similar to that of Theorem~1.1 in~\cite{MR4466386}. The latter gave an exponential inequality in the spirit of those of the paper for $U$-statistics, that is, sum of terms of the form $h\pr{X_1,\,\dots,\,X_r}$ where $\pr{X_i}_{i\,\geq\,1}$ is i.i.d.\ The connection with orthomartingales is the following. First, using decoupling (see~\cite{MR1334173}), the following inequality takes place:
\begin{multline}
\PP\left\{\left|\sum_{1\,\leq\,i_1\,<\,\dots\,i_r\,\leq\,n} h\left(X_{i_1},\,\dots,\,X_{i_r}\right) \right|>x\right\}\\
\leq C\PP\left\{\left|\sum_{1\,\leq\, i_1\,<\,\dots\,i_r\,\leq\,n} h\left(X_{i_1}^{\pr{1}},\,\dots,\,X_{i_r}^{\pr{r}}\right) \right|>Cx\right\},
\end{multline}
where $C$ is a universal constant and the sequences $\pr{X_i^{\pr{k}}}_{i\,\in\,\Z}, 1\leq k\leq r$ are mutually independent copies of $\pr{X_i}_{i\,\in\,\Z}$. If we assume that
\[
\bbE\left[h\left(X_1,\,\dots,\,X_r\right)\mid \sigma\left(X_i, i\in\left\{1,\,\dots,\,r\right\}\setminus \ens{k}\right)\right]=0
\]
for all $k\in \ens{1,\,\dots,\,r}$, then the random field $\pr{ h\pr{X_{i_1}^{\pr{1}},\,\dots,\,X_{i_r}^{\pr{r}}}}_{ i_1,\,\dots,\,i_r\,\in\,\Z}$ is an orthomartingale difference random field for the filtration $\pr{\Fca_{i_1,\,\dots,\,i_r}}_{i_1,\,\dots,\,i_r\,\in\,\Z}$ given by $\Fca_{i_1,\,\dots,\,i_r}=\sigma\pr{\eps_{j_q}^{\pr{q}}, j_q\leq i_q}$.

Let us now go into the details of the proof. Let $x$, $y>0$. We can assume without loss of generality that $x/y>3^d$. Indeed, suppose that we showed the existence of constants $A'_d$, $B_d$ and $C_d$ such that the inequality
\begin{multline}
\PP\left\{\max_{\gr{1}\,\imd\,\gri\,\imd\,\gr{n} } \left|S_{\gri} \right| >x\left|\grn\right|^{1/2}\right\}\leq A'_d\exp\left(-\left(\frac{x}{y}\right)^{2/d}\right)
\\+B_d\int_1^{+\infty}\PP\left\{\left|X_{\gr{1}}\right|>yuC_d\right\} u\left(\log\pr{1+u}\right)^{p_d}du,
\end{multline}
holds for each orthomartingale difference random field $\pr{X_{\gri}}_{\gri\,\in\,\Z^d}$ satisfying~\eqref{eq:invariance_loi_sommes} and for each $x,y>0$ such that $x/y>3^d$, then replacing $A'_d$ by $A_d=\max\ens{\exp\pr{9},A'_d}$ ensures that $A_d\exp\pr{-\pr{x/y}^{2/d}}
\geq 1$ when $x/y \leq 3^d$, so that the desired inequality becomes trivial in this range of parameters.
\begin{itemize}
\item[\eqref{step1}] Let $M$ and $M_N$ be defined as $M:=\max_{\gr{1}\,\imd\,\gri\,\imd\,\gr{n}} \abs{S_{\gri} }$, $\grn=\pr{n_1,\,\dots,\,n_d}$ and
\begin{equation}
M_N:=\max_{1\,\leq\,i_q\,\leq\,n_q,\,1\,\leq\,q\,\leq\,d-1}
\left|S_{i_1,\,\dots,\,i_{d-1},N}\right|.
\end{equation}
Define the filtration $\Gca_N$ as $\Fca_{n_1,\,\dots,\,n_{d-1}, N}$. We check that $\pr{M_N}_{N\,\geq\,1}$ is a submartingale with respect the filtration $\pr{\Gca_N}_{N\,\geq\,1}$. Indeed,
\begin{equation}
\bbE\left[M_N\middle| \Gca_{N-1}\right]
\geq \max_{1\,\leq\,i_q\,\leq\,n_q,\,1\,\leq\,q\,\leq\,d-1}
\left|\bbE\left[S_{i_1,\,\dots,\,i_{d-1},N}\right]\middle|\Gca_{N-1} \right|
\end{equation}
and by the orthomartingale property,
\begin{equation}
\bbE\left[S_{i_1,\,\dots,\,i_{d-1},N}\middle|\Gca_{N-1}\right] =S_{i_1,\,\dots,\,i_{d-1},N-1}.
\end{equation}
By Doob's inequality, we derive that
\begin{equation}
x\PP\ens{M>x}\leq \bbE\left[M_{n_d}\bfun\ens{M>x}\right].
\end{equation}
Expressing the latter expectation as an integral of the tail and cutting this integral at $x/2$ gives the bound
\begin{multline}\label{eq:first_step_induction}
\PP\left\{\max_{\gr{1}\,\imd\,\gri\,\imd\,\gr{n}} \left|S_{\gri} \right| >x\left|\grn\right|^{1/2}\right\} \\
\leq\int_1^{+\infty}
\PP\left\{
\max_{1\,\leq\,i_q\,\leq\,n_q,\,1\,\leq \,q\,\leq\,d-1}
\left|S_{i_1,\,\dots,\,i_{d-1},\,n_d}\right|>x\left|\grn\right|^{1/2}u/2 \right\}du.
\end{multline}
\item[\eqref{step2}] Define for $i_1,\,\dots,\,i_{d-1}$ the random variable
\begin{equation}
X'_{i_1,\,\dots,\,i_{d-1}}:=
\sum_{i_d=1}^{n_d}X_{i_1,\,\dots,\,i_{d-1},i_d}.
\end{equation}
Then $\pr{X'_{i_1,\,\dots,\,i_{d-1}}}_{i_1,\,\dots,\,i_{d-1}\,\in\,\Z}$ is an orthomartingale difference random field with respect to the commuting filtration $\pr{\Fca_{i_1,\,\dots,\,i_{d-1},n_d}}_{i_1,\,\dots,\,i_{d-1}\,\in\,\Z}$ satisfying~\eqref{eq:invariance_loi_sommes}. We thus apply the induction hypothesis with
\begin{equation}
\widetilde{x}:= xn_d^{1/2}u/2, \quad \widetilde{y}:= n_d^{1/2}x^{ 1/d}y^{\frac{d-1}d}u
\left(1+2\ln\pr{ u}\right)^{-\frac{d-1}2}/2.
\end{equation}
We get in view of~\eqref{eq:first_step_induction} and
\begin{equation}
\left(\frac{\widetilde{x}}{\widetilde{y}}\right)^{\frac{2}{d-1}}=
\left(\frac xy\right)^{2/d}\left(1+2\ln u\right)>3
\end{equation}
that
\begin{multline}\label{eq:second_step_induction}
\PP\left\{\max_{\gr{1}\,\imd\,\gri\,\imd\,\gr{n}} \left|S_{\gri}\right| >x\left(\grn\right)^{1/2}\right\} \leq A_{d-1}\int_1^{+\infty}\exp\left(-\left(\frac xy\right)^{2/d}\left(1+2\ln u\right)\right)du\\+ B_{d-1}\int_1^{+\infty}\int_1^{+\infty}\PP\left\{
\left|X'_{1,\,\dots,\,1}\right|>\frac{C_{d-1}}2vn_d^{1/2}x^{1/d}y^{\frac{d-1}d}u
\left(1+2\ln\pr{ u}\right)^{-\frac{d-1}2}\right\}\\
v\left(\log\pr{1+v}\right)^{p_{d-1}}dudv.
\end{multline}
For the first term, using that $ \pr{\frac xy}^{2/d}>3$, we get
\begin{multline}
\int_1^{+\infty}\exp\left(-\left(\frac xy\right)^{2/d}\left(1+2\ln u\right)\right)du
\\
\leq \exp\left(-\left(\frac xy\right)^{2/d}\right)
\int_1^{+\infty}\exp\left(-6\ln u\right)du
\end{multline}
hence
\begin{multline}\label{eq:second_step_induction_final}
\PP\left\{\max_{\gr{1}\,\imd\,\gri\,\imd\,\gr{n}} \left|S_{\gri}\right| >x\left|\grn\right|^{1/2}\right\} \leq
\frac{A_{d-1}}5\exp\left(-\left(\frac xy\right)^{2/d}\right)\\
+ B_{d-1}\int_1^{+\infty}\int_1^{+\infty}\PP\left\{
\left|\sum_{i_d=1}^{n_d}X_{1,\,\dots,\,1,i_d}\right|>\frac{C_{d-1}}2vn_d^{1/2}x^{1/d}y^{\frac{d-1}d}u
\left(1+2\ln\pr{ u}\right)^{-\frac{d-1}2} \right\}\\
v\left(\log\pr{1+v}\right)^{p_{d-1}}dudv.
\end{multline}
\item[\eqref{step3}] It remains to find a bound for the double integral. For fixed $u$ and $v>1$, we apply Proposition~\ref{prop:inegalite_deviation_martingales_dominee} in the following setting: $x$ is replaced by $\bar{x}$ defined by
\begin{equation}
\bar{x}:=\frac{C_{d-1}}2v x^{1/d}y^{\frac{d-1}d}u
\left(1+2\ln\pr{ u}\right)^{-\frac{d-1}2}
\end{equation}
and $y$ by $\bar{y}$ defined by
\begin{equation}
\bar{y}:= \frac{C_{d-1}}2yuv\left(1+2\ln u\right)^{-\frac{d-1}2} 2^{-1/2}\left(1+2\ln\left(uv\right)\right)^{-1/2}.
\end{equation}
In this way,
\begin{equation}
\frac 12\left(\frac{\bar{x}}{\bar{y}}\right)^2 =\left(\frac{x}y\right)^{2/d}\left(1+2\ln\left(uv\right)\right)
\end{equation}
and~\eqref{eq:inegalite_max_martingales_moments_ordre_p_domination_sto} gives
\begin{multline}\label{eq:estimated_last_step}
\PP\left\{\max_{\gr{1}\,\imd\,\gri\,\imd\,\gr{n}} \left|S_{\gri} \right| >x\left|\grn\right|^{1/2}\right\} \leq
\frac{A_{d-1}}5\exp\left(-\left(\frac xy\right)^{2/d}\right)\\+ 2B_{d-1}\exp\left(-\left(\frac{x}{y}\right)^{2/d}\right)\int_1^{+\infty}\int_1^{+\infty}
\exp\pr{-2\ln\pr{uv}}dudv\\
+ 4B_{d-1}\int_1^{+\infty}\int_1^{+\infty}\int_1^{+\infty} h\pr{u,v,w}dudvdw,
\end{multline}
where
\begin{multline}
h\pr{u,v,w}\\
{:=}\PP\left\{\left|X_{\gr{1}}\right|>
\frac{C_{d-1}}4yuvw\left(1+2\ln u\right)^{-\frac{d-1}2} 2^{-1/2}\left(1+2\ln\pr{uv}\right)^{-1/2} \right\}\\ vw\left(\log\pr{1+v}\right)^{p_{d-1}}.
\end{multline}
Observing that for $u,v>1$, $\pr{1+2\ln\pr{uv}}^{-1/2}\geq
\pr{1+2\ln\pr{u}}^{-1/2}\linebreak\pr{1+2\ln\pr{v}}^{-1/2}$ and defining the functions
\begin{equation}
f_q\pr{t}=t\left(1+2\ln t\right)^{-q}, t\geq 1, q>0
\end{equation}
\begin{equation}
g\colon t\mapsto \PP\left\{\left|X_{\gr{1}}\right|>
\frac{C_{d-1}}{4\cdot 2^{1/2}}yt \right\},
\end{equation}
we derive from estimate~\eqref{eq:estimated_last_step} that
\begin{multline}\label{eq:estimated_last_step_bis}
\PP\left\{\max_{\gr{1}\,\imd\,\gri\,\imd\,\gr{n}} \left|S_{\gri}\right| >x\left|\grn\right|^{1/2}\right\}\\
\leq
\frac{A_{d-1}}5\exp\left(-\left(\frac xy\right)^{2/d}\right)+ 2B_{d-1}\exp\left(-\left(\frac{x}{y}\right)^{2/d}\right) \\
+ 4B_{d-1}\int_1^{+\infty}\int_1^{+\infty}\int_1^{+\infty} g\left(w f_{d/2}\pr{u}f_{1/2}\pr{v}\right) vw\left(\log\pr{1+v}\right)^{p_{d-1}}dwdudv.
\end{multline}
Doing for fixed $u,v\geq 1$ the substitution $t=w f_{d/2}\pr{u}f_{1/2}\pr{v}$, we are reduced to show that there exists a constant $K$ such that for each $t\geq 1$,
\begin{equation}
I\pr{t}\leq K\left( \log\pr{1+t} \right)^{p_d},
\end{equation}
where
\begin{equation}
I\pr{t}:= \int_1^{+\infty}\int_1^{+\infty}\frac{v\left(\log\pr{1+v}\right)^{p_{d-1}} }{\left(f_{d/2}\pr{u}f_{1/2}\pr{v}\right)^2}
\bfun_{t\,>\,f_{d/2}\pr{u}f_{1/2}\pr{v}}du dv.
\end{equation}


Using the fact that there exists a constant $c$ such that for each $u\geq 1$, $f_{d/2}\pr{u}\geq c$, we derive that
\begin{equation}
I\pr{t}\leq \int_1^{+\infty}\int_1^{+\infty}\frac{v\left(\log\pr{1+v}\right)^{p_{d-1}} }{\left(f_{d/2}\pr{u}f_{1/2}\pr{v}\right)^2}
\bfun_{t\,>\,cf_{1/2}\pr{v} }du dv
\end{equation}
and since the integral over $u$ is convergent, we are reduced to show that there exists a constant $K$ such that for each $t\geq 1$,
\begin{equation}
\int_1^{+\infty}\frac1v\left(\log\pr{1+2v}\right)^{p_{d-1}+1}\bfun_{t\,\,>cf_{1/2}
\pr{v}}dv \leq K\left(\log\pr{1+t} \right)^{p_d}.
\end{equation}
Since there exists a constant $\kappa$ such that for each $v\geq 1$, $f_{1/2}\pr{v}\geq \kappa\sqrt v$, one can see that
\begin{multline}
\int_1^{+\infty}\frac1v\left(\log\pr{1+2v}\right)^{p_{d-1}+1}\bfun_{t\,>\,cf_{1/2}\pr{v} }dv \\
\leq 1_{t\,>\,c\kappa}\int_1^{\pr{t/(c\kappa)}^2}
\frac1v\left(\log\pr{1+2v}\right)^{p_{d-1}+1}dv
\end{multline}
and the fact that $p_d=p_{d-1}+2$ ends the proof of Theorem~\ref{thm:deviation_inequality_orthomartingale}.\qedhere
\end{itemize}
\end{proof}

\subsection{Proof of Theorem~\ref{thm:large_deviation}}
\begin{proof}
We apply Theorem~\ref{thm:deviation_inequality_orthomartingale} to $y=\abs{\gr{N}}^{1/\pr{2+d\gamma}}x^{2/\pr{2+d\gamma}}$. Bounding the resulting integral term by a constant times $\exp\pr{-y^\gamma}$, one sees that in this case, the two terms of the right hand side of~\eqref{eq:exp_deviation_inequality_orthomartingale} have a similar contribution.
\end{proof}

\subsection{Proof of Proposition~\ref{prop:tightness_partial_sums}}

\begin{proof}
The proof of the tightness criterion rests on the Schauder decomposition of the spaces $\clH_\rho^o\left([0,1]^d\right)$. In order to state it, we need to introduce the following notations.

Set for $j\geqslant 0$,
\begin{equation}
W_j:=\left\{k2^{-j},0\leq k\leq 2^j\right\}^d
\end{equation}
and
\begin{equation}
V_0:=W_0, \quad V_j:=W_j\setminus W_{j-1},j\geqslant 1.
\end{equation}

We define for $\gr{v}\in V_j$ the pyramidal function $\Lambda_{j,
\gr{v}}$ by
\begin{equation}
\Lambda_{j,\gr{v}}(\gr{t}):=\Lambda\left(2^j(\gr{t}-\gr{v})\right), \quad
\gr{t}\in [0,1]^d,
\end{equation}
where
\begin{equation}
\Lambda(\gr{t}):=\max\left\{0,1-\max_{t_i\,<\,0}\abs{t_i}-\max_{t_i\,>\,0}\abs{t_i}\right\},
\quad \gr{t}=(t_i)_{i=1}^d\in [-1,1]^d.
\end{equation}
For $x\in \clH_\alpha^o([0,1]^d)$, we define the coefficients $\lambda_{j,\gr{v}}(x)$ by $\lambda_{0,\gr{v}}(x)=x(\gr{v})$, $\gr{v}\in V_0$ and for $j\geqslant 1$ and $v\in V_j$,
\begin{equation}
\lambda_{j,\gr{v}}(x):=x\left(\gr{v}\right)-\frac 12\left(x\left(\gr{v^-}\right)+ x\left(\gr{v^+}\right)\right),
\end{equation}
where $\gr{v^+}$ and $\gr{v^-}$ are define in the following way. Each $\gr{v}\in V_j$ is represented in a unique way by $\gr{v}=
(k_i2^{-j})_{i=1}^d$. Then $\gr{v^+}:=(v^+_i)_{i=1}^d$ and $\gr{v^-}:=(v^-_i)_{i=1}^d$ are defined by
\begin{equation}
\gr{v_i^-}:=
\begin{cases}
v_i-2^{-j},&\text{ if }k_i\text{ is odd;}\\
v_i,&\text{ if }k_i\text{ is even}
\end{cases}\quad
\gr{v_i^+}:=
\begin{cases}
v_i+2^{-j},&\text{ if }k_i\text{ is odd;}\\
v_i,&\text{ if }k_i\text{ is even.}
\end{cases}
\end{equation}
The sequential norm is defined by
\begin{equation}
\norm{x}_\rho^{\seq}:=
\sup_{j\,\geqslant\,0}\rho\left(2^{-j}\right)^{-1}\max_{\gr{v}\,\in\,V_j}
\left|\lambda_{j,\gr{v}}(x)\right|, \quad x\in \clH_\rho^o\left([0,1]^d\right).
\end{equation}
By~\cite{MR2141332}, the norm $\norm{\cdot}_\rho^{\seq}$ is equivalent to $\norm{\cdot}_\rho$ on $\clH_\rho^o([0,1]^d)$.

A general tightness criterion is available for moduli of the form $\rho\colon h\mapsto h^\alpha$. The criterion rests on a Schauder decomposition of $\clH_\rho^o$ as $\bigoplus_{j\,\geq\,1}E_j$, where $E_j$ is the vector space generated by the functions $\Lambda_{j,\gr{v}}$, $\gr{v}\in V_j$.

\begin{theo}[{\cite[Theorem~6]{MR2340883}}]\label{thm:general_theorem_tightness_H\"older_dim_d}
Let $\ens{\zeta_{\grn},\grn\in \N^d}$ and $\zeta$ be random elements with values in the space $\clH_\alpha([0,1]^d)$. Assume that the following conditions are satisfied.
\begin{enumerate}
\item For each dyadic $\gr{t}\in [0,1]^d$, the net $\ens{\zeta_{\grn}(\gr{t}),
\grn\in \N^d}$ is asymptotically tight on $\R$.
\item For each positive $\varepsilon$,
\begin{equation}
\lim_{J\,\to\,\infty}\limsup_{\min\,\grn\,\to\,\infty}
\PP\left\{\sup_{j\,\geqslant\,J}2^{\alpha j}\max_{\gr{v}\,\in\,V_j}\left|\lambda_{j,
\gr{v}}(\zeta_{\grn}) \right|>\varepsilon\right\}=0.
\end{equation}
\end{enumerate}
Then the net $\ens{\zeta_{\grn},\,\grn\,\in\,\N^d}$ is asymptotically tight in the space $\clH_\alpha([0,1]^d)$.
\end{theo}

This extends readily to $\rho\in\Rca_{1/2,d}$.

First observe that by bounding from below the sum over $j$ by the term at index $j=J$ and taking $i=d$, \eqref{eq:tightness_random_fields_sufficient_cond} implies that
\begin{equation}
\lim_{J\to \infty}\limsup_{\min\,\grm\,\to\,\infty} 2^J \PP\left\{\max_{\gr{1}\,\imd\,\gr{k}\,\imd\,2^{\grm -J\gr{e}_d } }
\left|\sum_{\gr{1}\,\imd\,\gri\,\imd\,\grk} X_{\gri} \right|> \varepsilon\rho\left(2^{-j}\right)
\prod_{u=1}^d2^{m_u/2}\right\}=0,
\end{equation}
and doing the replacement of index $m'_d=m_d-J$ for a fixed $J$ gives asymptotic tightness of $ \pr{W_{\grn}\pr{\gr{t}}}_{\grn\,\smd\,\gr{1}}$ for each $\gr{t}$.

It remains to check the second condition of Theorem~\ref{thm:general_theorem_tightness_H\"older_dim_d}. Since the Schauder decomposition $\clH_\rho^o=
\bigoplus_{j\,\geq\,1}E_j$ is also valid for $\rho\in\Rca_{1/2,d}$, this theorem also holds with the map $h\mapsto h^\alpha$ replaced by $\rho$. Therefore, it suffices to prove that
\begin{equation}\label{eq:step_for_tightness}
\lim_{J\,\to\,\infty}\limsup_{\min\,\grn\,\to\,\infty}
\PP\left\{\sup_{j\,\geqslant\,J} \rho\left(2^{-j}\right)^{-1}\max_{\gr{v}\,\in\,V_j}\left|\lambda_{j,
\gr{v}}\left(W_{\grn}\right) \right|>\varepsilon \right\}=0.
\end{equation}


Consider for $\gr{s}:=(s_2,\dots,s_d)\in [0,1]^{d-1}$ and $t,t'\in [0,1]$ the quantity
\begin{equation}
\Delta_{\grn}(t,t',\gr{s}):=\sqrt{\left|\grn\right|}
\,\left|W_{\grn}\left(t',\gr{s}\right)-W_{\grn}\left(t,\gr{s}\right)\right|.
\end{equation}

We recall the following lemma:
\begin{lemm}[{\cite[Lemma~11]{MR2340883}}]\label{lem:Rac_Suq_Zem}
For any $t',t\in [0,1]$, $t'>t$, the following inequality holds:
\begin{multline}
\sup_{\gr{s}\,\in\,[0,1]^{d-1}}\Delta_{\grn}\left(t,t',\gr{s}\right)
\leq 3^d\bfun\left\{t'-t\geqslant \frac 1{n_1}\right\}\quad
\max_{\mathclap{\substack{1\,\leq\,k_\ell\,\leq\,n_\ell \\
2\,\leq\,\ell\,\leq\,d} }}\quad\left|\sum_{i_1=[n_1t]+1}^{[n_1t']}\quad
\sum_{\mathclap{\substack{1\,\leq\,i_\ell\,\leq\,k_\ell \\
2\,\leq\,\ell\,\leq\,d} }}\quad X_{\gri}\right|\\
+3^d\min\left\{1,n_1(t'-t)\right\}\max_{1\,\leq\,i_1\,\leq\,n_1}
\quad
\max_{\mathclap{\substack{1\,\leq\,k_\ell\,\leq\,n_\ell \\
2\,\leq\,\ell\,\leq\,d} }}\quad\left|\quad
\sum_{\mathclap{\substack{1\,\leq\,i_\ell\,\leq\,k_\ell \\
2\,\leq\,\ell\,\leq\,d}}}\quad X_{\gri}\right|.
\end{multline}
\end{lemm}

Now, we define for $q\in [d]$ and $\gr{s}=\pr{s_\ell}_{\ell\in [d] \setminus \ens q}\in [0,1]^{d-1}$,
\begin{multline*}
\Delta_{\grn}^{(q)}(t,t',\gr{s}):=\\
\left|W_{\grn}\left(s_1,\,\dots,\,s_{q-1},t',s_{q+1},\,\dots,\,s_d\right)- W_{\grn}\left(s_1,\,\dots,\,s_{q-1},t,s_{q+1},\,\dots,\,s_d\right)\right|.
\end{multline*}
By definition of $\lambda_{j,\gr{v}}$ and $V_j$, the inequality
\begin{multline}\label{eq:H\"older_modulus_random_fields_probability_bound_proof}
\PP\left\{\sup_{j\,\geqslant\,J}\rho\left(2^{-j}\right)^{-1}\max_{\gr{v}\,\in\,V_j}\left|\lambda_{j,\gr{v}}\left(W_{\grn}\right) \right|>2^{d+1} \eps\right\} \\ \leq
\sum_{q=1}^d\PP\left\{\sup_{j\,\geqslant\,J}\rho\left(2^{-j}\right)^{-1}
\max_{\mathclap{\substack{0\,\leq\,k\,<\,2^j\\
\gr{0}\,\imd\,\gru\,\imd\,2^j\,\gr{1}}}} \Delta^{(q)}_{\grn}\left(t_{k+1},t_k;
\gr{s_u}\right)>2\cdot \eps\right\}
\end{multline}
takes place, where $t_k=k2^{-j}$ and $\gr{s_u}=(u_i2^{-j})_{i\,\in\,[d] \setminus\ens{q}}$. We show that for each positive~$\eps$,

\begin{equation}
\lim_{J\,\to\,\infty}\limsup_{\grn\,\to\,\infty}
\PP\left\{\sup_{j\,\geqslant\,J}\rho\left(2^{-j}\right)^{-1}
\max_{\mathclap{\substack{0\,\leq\,k\,<\,2^j\\
\gr{0}\,\imd\,\gru\,\imd\,2^j\gr{1}}}} \Delta^{(d)}_{\grn}\left(t_{k+1},t_k;
\gr{s_u}\right)>2\cdot 3^d\cdot \eps\right\},
\end{equation}
the treatment of the corresponding terms with $\Delta^{(q)}$ instead of $\Delta^{(d)}$ can be done by switching the roles of the coordinates. We will use the following notations: we will denote by $\gri,\grk,\grn,\grm$ elements of $\Z^d$ and $\gr{i'},\gr{k'},\gr{n'},\gr{m'}$ elements of $\Z^{d-1}$ and $\pr{\pr{\gr{i'},i_d}}$ will denote the element of $\Z^d$ whose first $d-1$ coordinates are those of $\gr{i'}$ and $d$th one is $i_d$, and similarly for other letters. We will denote by $\imd$ the coordinatewise order on $\Z^d$ and $\Z^{d-1}$, since there will be no ambiguity. Finally, we define let $\gr{1'}=\sum_{q=1}^{d-1}\gr{e}_q$.

We have in view of Lemma~\ref{lem:Rac_Suq_Zem} that
\begin{multline}\label{eq:H\"older_modulus_random_fields_probability_bound_proof_2}
\PP\left\{\sup_{j\,\geqslant\,J}\rho\left(2^{-j}\right)^{-1}
\max_{\mathclap{\substack{0\,\leq\,k<2^j\\
\gr{0}\,\imd\,\gru\,\imd\,2^j\gr{1}}}}\Delta^{\pr{d}}_{\grn}\left(t_{k+1},t_k;
\gr{s_u}\right)>2\cdot 3^d\eps\left|\grn\right|^{1/2}\right\}\\
\begin{aligned}
&\leq \PP\left\{\sup_{j\,\geqslant\,J}
\max_{0\,\leq\,a\,<\,2^j}\rho\left(2^{-j}\right)^{-1} \bfun\left\{2^{-j}\geqslant \frac 1{n_d}\right\}
\max_{\gr{1'}\,\imd\,\gr{k'}\,\imd\,\gr{n'}}
\left|\sum_{i_d=\left[n_da2^{-j}\right]+1}^{\left[n_d(a+1)2^{-j}\right]}
\sum_{\gr{1'}\,\imd\,\gr{i'}\,\imd\,\gr{k'}} X_{\left(\gr{i'},i_d\right) } \right|\right.\\
&+\left.\min\left\{1,n_d2^{-j}\right\}\max_{1\,\leq\, i_d\,\leq\,n_d}
\max_{\gr{1'}\,\imd\,\gr{k'}\,\imd\,\gr{n'}} \left|
\sum_{\gr{1'}\,\imd\,\gr{i'}\,\imd\,\gr{k'}} X_{\left(\gr{i'},i_d\right) }\right| >2\eps\left|\grn\right|^{1/2}\right\} \\
&\leq \PP\left\{\sup_{j\,\geqslant\,J}\max_{0\,\leq \,a\,<\,2^j}\rho\left(2^{-j}\right)^{-1}
\bfun\left\{2^{-j}\geqslant \frac 1{n_d}\right\}\right.
\\
&\mkern150mu\left.\max_{\gr{1'}\,\imd\,\gr{k'}\,\imd\,\gr{n'}}
\left|\sum_{i_d=\left[n_da2^{-j}\right]+1}^{\left[n_d(a+1)2^{-j}\right]}
\sum_{\gr{1'}\,\imd\,\gr{i'}\,\imd\,\gr{k'}} X_{\left(\gr{i'},i_d\right) } \right|>\eps\abs{\grn}^{1/2}\right\}\\
&+\PP\left\{\sup_{j\,\geqslant\,J}\rho\left(2^{-j}\right)^{-1}
\min\left\{1,n_d2^{-j}\right\}\max_{1\,\leq\,i_d\,\leq \,n_d}
\max_{\gr{1'}\,\imd\,\gr{k'}\,\imd\,\gr{n'}} \left|
\sum_{\gr{1'}\,\imd\,\gr{i'}\,\imd\,\gr{k'} } X_{\left(\gr{i'},i_d\right) }\right|>\eps\left|\grn\right|^{1/2}\right\}.
\end{aligned}
\end{multline}
Since the indicator in the first term of the right hand side of~\eqref{eq:H\"older_modulus_random_fields_probability_bound_proof_2} vanishes if $j>\log n_d$, we have
\begingroup
\allowdisplaybreaks
\begin{align*}
&\PP\left\{\sup_{j\,\geqslant\,J}\max_{0\,\leq\,a\,<\,2^j}\rho\left(2^{-j}\right)^{-1}
\bfun\left\{2^{-j}\geqslant \frac 1{n_d}\right\}\right.
\\
&\mkern100mu\left. \max_{\gr{1'}\,\imd\,\gr{k'}\,\imd\,\gr{n'}}
\left|\sum_{i_d=\left[n_da2^{-j}\right]+1}^{\left[n_d(a+1)2^{-j}\right]}
\sum_{\gr{1'}\,\imd\,\gr{i'}\,\imd\,\gr{k'}} X_{\left(\gr{i'},i_d\right) } \right|>\eps\left|\grn\right|^{1/2}\right\}\\
\leq &\PP\left\{\sup_{J\,\leq\,j\,\leq\,\log\,n_d}\max_{0\,\leq\,a\,<\,2^j}\rho\left(2^{-j}\right)^{-1}
\max_{\gr{1'}\,\imd\,\gr{k'}\,\imd\,\gr{n'}}
\left|\sum_{i_d=\left[n_da2^{-j}\right]+1}^{\left[n_d(a+1)2^{-j}\right]}
\sum_{\gr{1'}\,\imd\,\gr{i'}\,\imd\,\gr{k'}} X_{\left(\gr{i'},i_d\right) }\right|>\eps\left|\grn\right|^{1/2}\right\}\\
\leq &\sum_{j=J}^{\log n_d}2^j
\max_{0\,\leq\,a\,<\,2^j}\PP\left\{
\rho\left(2^{-j}\right)^{-1}
\max_{\gr{1'}\,\imd\,\gr{k'}\,\imd\,\gr{n'}}
\left|\sum_{i_d=\left[n_da2^{-j}\right]+1}^{\left[n_d(a+1)2^{-j}\right]}
\sum_{\gr{1'}\imd\gr{i'}\imd\gr{k'} } X_{\left(\gr{i'},i_d\right) } \right|>\eps\abs{\grn}^{1/2} \right\},
\end{align*}
\endgroup
and by stationarity, it follows that
\begin{multline*}
\PP\left\{\sup_{j\,\geqslant\,J}\max_{0\,\leq\,a\,<\,2^j}\rho\left(2^{-j}\right)^{-1}
\bfun\left\{2^{-j}\geqslant \frac 1{n_d}\right\}\right. \\
\left.\max_{\gr{1'}\,\imd\,
\gr{k'}\,\imd\,\gr{n'} }
\left|\sum_{i_d=\left[n_da2^{-j}\right]+1}^{\left[n_d(a+1)2^{-j}\right]}
\sum_{\gr{1'}\,\imd\,\gr{i'}\,\imd\,\gr{k'}} X_{\left(\gr{i'},i_d\right) }\right|>\eps\abs{\grn}^{1/2}\right\}\\
\leq \sum_{j=J}^{\log n_d}2^j
\PP\left\{\rho\left(2^{-j}\right)^{-1}
\max_{1\,\leq\,k_d\,\leq\,2n_d2^{-j}} \max_{\gr{1'}\,\imd\,\gr{k'}\,\imd\,\gr{n'}}
\left|\sum_{i_d= 1}^{k_d}
\sum_{\gr{1'}\,\imd\,\gr{i'}\,\imd\,\gr{k'}} X_{\left(\gr{i'},i_d\right) } \right|>\eps\abs{\grn}^{1/2} \right\}.
\end{multline*}
If $\grn =\pr{n_1,\,\dots,\,n_d}$ is such that $2^{m_i}\leq n_i\leq 2^{m_i+1}-1$ for each $i\in [d]$, then we derive that
\begin{multline}\label{eq:H\"older_modulus_random_fields_probability_bound_proof_3}
\PP\left\{\sup_{j\,\geqslant\,J}\max_{0\,\leq\,a\,<\,2^j}\rho\left(2^{-j}\right)^{-1}
\bfun\left\{2^{-j}\geqslant \frac 1{n_d}\right\}\right.
\\
\left. \max_{\gr{1'}\,\imd\,
\gr{k'}\,\imd\,\gr{n'} }
\left|\sum_{i_d=\left[n_da2^{-j}\right]+1}^{\left[n_d(a+1)2^{-j}\right]}
\sum_{\gr{1'}\,\imd\,\gr{i'}\,\imd\,\gr{k'}} X_{\left(\gr{i'},i_d\right) } \right|>\eps\abs{\grn}^{1/2}\right\}\\
\leq \sum_{j=J}^{m_d+1}2^j
\PP\left\{\rho\left(2^{-j}\right)^{-1}
\max_{\gr{1}\,\imd\,\gr{k}\,\imd\,\gr{2^{\grm-
\pr{j-2}\gr{e_d}}} }
\left|
\sum_{\gr{1}\,\imd\,\gr{i}\,\imd\,\gr{k}} X_{\left(\gr{i'},i_d\right) } \right|>\eps \prod_{u=1}^d 2^{m_u/2} \right\}.
\end{multline}
For the second term of the right hand side of~\eqref{eq:H\"older_modulus_random_fields_probability_bound_proof_2}, notice that
\begin{equation}
\sup_{j\,\geqslant\,J}\rho\left(2^{-j}\right)^{-1}\min\left\{1,n_d2^{-j}\right\}\leq c_\rho
\rho\pr{1/n_d}^{-1}.
\end{equation}
Indeed, if $j\leq \log n_d$, then $2^j\leq n_d$ hence $\rho\pr{2^{-j}}^{-1}\min\{1,n_d2^{-j}\}=\rho\pr{2^{-j}}^{-1}$ and the sequence $\pr{\rho\pr{2^{-j}}^{-1}}_j$ is increasing, and if $j>\log n_d$, then $2^j>n_d$, hence $\min\{1,n_d2^{-j}\}=n_d2^{-j}$ and for such $j$'s, we have $\rho\pr{2^{-j}}^{-1}n_d2^{-j}$ and we use decreasingness of the sequence $\pr{\rho\pr{2^{-j}}^{-1}2^{-j}}_j$. As a consequence, after having bounded the probability of the max over $i_d$ by the sum of probabilities and used stationarity, we obtain
\begin{multline}
\PP\left\{\sup_{j\,\geqslant\,J}\rho\left(2^{-j}\right)^{-1}
\min\left\{1,n_d2^{-j}\right\}\max_{1\,\leq\,i_d\,\leq\, n_d}\right.\\
\left.\max_{\gr{1'}\,\imd\,\gr{k'}\,\imd\,\gr{n'}} \left|
\sum_{\gr{1'}\,\imd\,\gr{i'}\,\imd\,\gr{k'}} X_{\left(\gr{i'},i_d\right) }\right|>\eps\abs{\grn}^{1/2}\right\}\\
\leq n_d
\PP\left\{ \max_{\gr{1'}\,\imd\,\gr{k'}\,\imd\,\gr{n'}} \left|
\sum_{\gr{1'}\,\imd\,\gr{i'}\,\imd\,\gr{k'}\,} X_{\left(\gr{i'},1\right)}\right|>\eps\abs{\grn}^{1/2}\rho\pr{1/n_d} /c_{\rho}\right\}.
\end{multline}
Notice that if $\grn =\pr{n_1,\,\dots,\,n_d}$ is such that $2^{m_i}\leq n_i\leq 2^{m_i+1}-1$ for each $i\in [d]$, then for each $j\leq m_d$,
\begin{multline*}
\frac{1}{\abs{\grn}^{1/2}\rho\pr{1/n_d}}\max_{\gr{1'}\,\,\imd\,\gr{k'}\,\imd\,\gr{n'}}
\left|
\sum_{\gr{1'}\,\imd\,\gr{i'}\,\imd\,\gr{k'}} X_{\left(\gr{i'},1\right)}\right|\\
\begin{aligned}
&\leq
\frac{1}{\rho\left(2^{-m_{d}-1}\right)\prod_{u=1}^d2^{m_u/2}}
\max_{\gr{1'}\,\imd\,\gr{k'}\,\imd
\,\gr{2^{m'}}}\left|\sum_{\gr{1'}\,\imd\,\gr{i'}\,\imd\,\gr{k'}} X_{\left(\gr{i'},1\right)}\right|\\
&\leq
\frac{1}{\rho\left(2^{-m_{d}-1}\right)\prod_{u=1}^d2^{m_u/2}}
\max_{1\,\leq\,k_d\,\leq\,2^{m_d+1-j}} \max_{\gr{1'}\,\imd\,\gr{k'}\,\imd\,
\gr{2^{m'}}}\left|\sum_{i_d=1}^{k_d}\sum_{\gr{1'}\,\imd\,\gr{i'}\,\imd\,\gr{k'}} X_{\left(\gr{i'},i_d\right)}\right|
\end{aligned}
\end{multline*}
hence
\begin{multline}\label{eq:H\"older_modulus_random_fields_probability_bound_proof_4}
\PP\left\{\sup_{j\,\geqslant\,J}\rho\left(2^{-j}\right)^{-1}
\min\left\{1,n_d2^{-j}\right\}\max_{1\,\leq\,i_d\,\leq \,n_d}\right.\\
\left.\max_{\gr{1'}\,\imd\,\gr{k'}\,\imd\,\gr{n'}} \left|
\sum_{\gr{1'}\,\imd\,\gr{i'}\,\imd\,\gr{k'}} X_{\left(\gr{i'},i_d\right) }\right|>\eps\abs{\grn}^{1/2}\right\}\\
\leq \sum_{j=J}^{m_d+1}2^j \PP\left\{
\max_{\gr{1}\,\imd\,\gr{k}\,\imd\,
\gr{2^{m-\pr{j-1}\gr{e_d}}}}\left|\sum_{\gr{1}\,\imd\,\gr{i}\,\imd\,\grk } X_{\gri}\right|> \eps\rho\left(2^{-j}\right)
\prod_{u=1}^d2^{m_u/2}/c_\rho\right\}.
\end{multline}

Combining inequalities~\eqref{eq:H\"older_modulus_random_fields_probability_bound_proof_2} with~\eqref{eq:H\"older_modulus_random_fields_probability_bound_proof_3} and~\eqref{eq:H\"older_modulus_random_fields_probability_bound_proof_4}, we obtain~\eqref{eq:step_for_tightness}. This ends the proof of Proposition~\ref{prop:tightness_partial_sums}.
\end{proof}

\subsection{Proof of Theorem~\ref{thm:WIP_Holder}}

We have to check that~\eqref{eq:tightness_random_fields_sufficient_cond} is satisfied. For simplicity, we will do this for $q=d$; the general case can be done similarly. To this aim, we apply inequality~\eqref{eq:exp_deviation_inequality_orthomartingale} for fixed $\gr{m}\smd\gr{1}$, $J\geq 1$ and $j\in\ens{J,\,\dots,\,m_d}$ in the following setting: $\widetilde{\gr{m}}=\pr{2^{m_1},\,\dots,\,2^{m_{d-1}},2^{m_d-j}}$, $x=\eps \rho\pr{2^{-j}}2^{j/2}$ and $y= \eps/2\rho\pr{2^{-j}}2^{j/2} j^{-d/2}$. The sum of the exponential terms in~\eqref{eq:exp_deviation_inequality_orthomartingale} can be bounded by the remainder of a convergent series. With the assumption on $\rho$, the sum of the obtained integral terms does not exceed $\sum_{j\,\geq\,J}2^j\int_1^{+\infty}\PP\ens{\abs{X_{\gr{1}}}>L\pr{2^j}uC} u\pr{\log\pr{1+u}}^{p_d}du$ where $C$ depends only on $\rho$ and~$\eps$.

Now, we will show that~\eqref{eq:cond_suff_WIP} guarantees the convergence to zero of the previous term as $J$ goes to infinity. As there exists a constant $\kappa$ such that $u\pr{\log\pr{1+u}}^{p_d}\leq \kappa u^2$ for each $u\geq 1$, it suffices to prove that for each $C$,
\begin{equation}
\sum_{j\,\geq\,1}2^j\int_1^{+\infty}\PP\left\{\left|X_{\gr{1}}\right|>L\left(2^j\right)uC\right\} u^{2}du<\infty
\end{equation}
This will be a consequence of the following lemmas.

\begin{lemm}\label{lem:slowly_varying_2j_sur_L2j}
Let $L\colon \R_{+}\to \R_{+}$ be a slowly varying function. There exists a constant $C_L$ such that for each $k\geq 1$,
\begin{equation}\label{eq:slowly_varying_2j_sur_L2j}
\sum_{j=1}^k\frac{2^j}{L\left(2^j\right)}\leq C_L\frac{2^k}{L\left(2^k\right)}.
\end{equation}
\end{lemm}

\begin{proof}
From Potter's bound (see Lemma~1.5.6. in~\cite{MR1015093}), there exists a constant $K$ such that for each $1\leq j\leq k$, $L\pr{2^k} \leq K L\pr{2^j} 2^{\frac{k-j}2}$. Consequently,
\begin{equation}
\sum_{j=1}^k\frac{2^j}{L\left(2^j\right)} =\sum_{j=1}^k\frac{2^j}{L\left(2^k\right)}\frac{L\left(2^k\right)}{L\left(2^j\right)}\leq K\sum_{j=1}^k\frac{2^j}{L\left(2^k\right)}2^{\frac{k-j}2}=K\frac{2^{k/2}}{L\left(2^k\right)}\sum_{j=1}^k2^{\frac j2}
\end{equation}
and the change of index $\ell=k-j$ shows that we can take $C_L=K\sum_{\ell\,\geq\,0}2^{-\ell/2}$. This ends the proof of Lemma~\ref{lem:slowly_varying_2j_sur_L2j}.
\end{proof}


\begin{lemm}\label{lem:slowly_varying}
Let $X$ be a non-negative random variable and let $L\colon \R_{+}\to \R_{+}$ be a slowly varying increasing function such that $L\pr{x}\to \infty$ as $x\to \infty$. Suppose that
\begin{equation}
\forall A>0, \sum_{j\,\geq\,1}2^j\PP\left\{X>L\left(2^j\right)A\right\}<\infty.
\end{equation}
Then for each $C>0$,
\begin{equation}
\sum_{j\,\geq\,1}2^j\int_1^{+\infty}\PP\left\{X>L\left(2^j\right)uC\right\} u^{2}du<\infty.
\end{equation}
\end{lemm}
\begin{proof}
Without loss of generality, we can assume that $C=1$. First, observe that
\begin{equation}
\int_1^{+\infty}\PP\left\{X>L\left(2^j\right)u\right\} u^{2}du\leq \left[\left(\frac{X}{L\left(2^j\right)}\right)^{3}\bfun\left\{X>L\left(2^j\right) \right\}\right].
\end{equation}
Therefore,
\begin{multline}
\sum_{j\,\geq\,1}2^j\int_1^{+\infty}\PP\left\{X>L\left(2^j\right)u\right\} u^{2}du \\
\leq \sum_{j\,\geq\,1}2^j\left[\left(\frac{X}{L\left(2^j\right)}\right)^{3}\sum_{k\,\geq\,j}\bfun\left\{ L\left(2^k\right)<X\leq L\left(2^{k+1}\right) \right\}\right]\\
\leq \sum_{j\,\geq\,1}2^j\left[\sum_{k\,\geq\,j}\left(\frac{L\left(2^{k+1}\right) }{L\left(2^j\right)}\right)^{3}\bfun\left\{ L\left(2^k\right)<X\leq L\left(2^{k+1}\right) \right\}\right].
\end{multline}
Switching the sums over $j$ and $k$ and using Lemma~\ref{lem:slowly_varying} with $L^3$ instead of $L$ reduces use to show that
\begin{equation}
\sum_{k\,\geq\,1} 2^k\PP\left\{ L\left(2^k\right)<X\leq L\left(2^{k+1}\right) \right\}<\infty,
\end{equation}
which follows from the assumption and the obvious bound $\PP\{ L\pr{2^k}<X\leq L\pr{2^{k+1}} \} \linebreak\leq
\PP\{ L\pr{2^k}<X \}$.
\end{proof}



The last part of the statement of Theorem~\ref{thm:WIP_Holder} follow from~\cite{MR3913270}. This ends the proof of Theorem~\ref{thm:WIP_Holder}.


\subsection*{Acknowledgement}
The author is grateful to the referee for many useful comments and corrections that improved the paper.


\begin{thebibliography}{99}
\bibitem{MR1015093} Bingham, Nicholas; Goldie, Charles; Teugels, Jan Regular variation, Encyclopedia of Mathematics and Its Applications, 27, Cambridge University Press, 1989 | Zbl
\bibitem{MR0126871} Billingsley, Patrick The Lindeberg–Lévy theorem for martingales, Proc. Am. Math. Soc., Volume 12 (1961), pp. 788-792 | Zbl
\bibitem{MR0254912} Cairoli, Renzo Un théorème de convergence pour martingales à indices multiples, C. R. Acad. Sci., Paris, Sér. A, Volume 269 (1969), p. A587-A589 | MR | Zbl
\bibitem{MR3504508} Cuny, Christophe; Dedecker, Jérôme; Volný, Dalibor A functional CLT for fields of commuting transformations via martingale approximation, Zap. Nauchn. Semin. (POMI), Volume 441 (2015) no. 22, pp. 239-262 | Zbl
\bibitem{MR1875665} Dedecker, Jérôme Exponential inequalities and functional central limit theorems for a random fields, ESAIM, Probab. Stat., Volume 5 (2001), pp. 77-104 | MR | DOI | Zbl
\bibitem{MR1334173} de la Peña, Victor H.; Montgomery-Smith, Stephen J. Decoupling inequalities for the tail probabilities of multivariate $U$-statistics, Ann. Probab., Volume 23 (1995) no. 2, pp. 806-816 | MR | Zbl
\bibitem{MR30714} Erdös, Pál On a theorem of Hsu and Robbins, Ann. Math., Volume 20 (1949), pp. 286-291 | MR | Zbl
\bibitem{MR3005732} Fan, Xiequan; Grama, Ion; Liu, Quansheng Large deviation exponential inequalities for supermartingales, Electron. Commun. Probab., Volume 17 (2012) no. 59, p. 8 | MR | Zbl
\bibitem{MR3311214} Fan, Xiequan; Grama, Ion; Liu, Quansheng Exponential inequalities for martingales with applications, Electron. J. Probab., Volume 20 (2015) no. 1, 1, 22 pages | MR | Zbl
\bibitem{MR3426520} Giraudo, Davide Holderian weak invariance principle under a Hannan type condition, Stochastic Processes Appl., Volume 126 (2016) no. 1, pp. 290-311 | MR | DOI | Zbl
\bibitem{MR3615086} Giraudo, Davide Holderian weak invariance principle for stationary mixing sequences, J. Theor. Probab., Volume 30 (2017) no. 1, pp. 196-211 | DOI | MR | Zbl
\bibitem{MR3770868} Giraudo, Davide Hölderian weak invariance principle under the Maxwell and Woodroofe condition, Braz. J. Probab. Stat., Volume 32 (2018) no. 1, pp. 172-187 | MR | Zbl
\bibitem{MR3869881} Giraudo, Davide Invariance principle via orthomartingale approximation, Stoch. Dyn., Volume 18 (2018) no. 6, 1850043, 29 pages | MR | Zbl
\bibitem{MR4046858} Giraudo, Davide Deviation inequalities for Banach space valued martingales differences sequences and random fields, ESAIM, Probab. Stat., Volume 23 (2019), pp. 922-946 | MR | DOI | Zbl
\bibitem{MR4466386} Giraudo, Davide An exponential inequality for $U$-statistics of i.i.d. data, Theory Probab. Appl., Volume 66 (2021) no. 3, pp. 408-429 | MR | DOI | Zbl
\bibitem{MR1759810} Hamadouche, Djamel Invariance principles in Hölder spaces, Port. Math., Volume 57 (2000) no. 2, pp. 127-151 | MR | Zbl
\bibitem{MR19852} Hsu, Pao Lu; Robbins, Herbert Complete convergence and the law of large numbers, Proc. Natl. Acad. Sci. USA, Volume 33 (1947), pp. 25-31 | Zbl | MR
\bibitem{MR0151997} Ibragimov, Il’dar A. A central limit theorem for a class of dependent random variables, Teor. Veroyatn. Primen., Volume 8 (1963), pp. 89-94 | Zbl | MR
\bibitem{MR2794415} Kuczmaszewska, Anna; Lagodowski, Zbigniew A. Convergence rates in the SLLN for some classes of dependent random fields, J. Math. Anal. Appl., Volume 380 (2011) no. 2, pp. 571-584 | Zbl | MR | DOI
\bibitem{MR3451971} Lagodowski, Zbigniew A. An approach to complete convergence theorems for dependent random fields via application of Fuk-Nagaev inequality, J. Math. Anal. Appl., Volume 437 (2016) no. 1, pp. 380-395 | Zbl | MR | DOI
\bibitem{MR1856684} Lesigne, Emmanuel; Volný, Dalibor Large deviations for martingales, Stochastic Processes Appl., Volume 96 (2001) no. 1, pp. 143-159 | Zbl | MR | DOI
\bibitem{MR2568267} Liu, Quansheng; Watbled, Frédérique Exponential inequalities for martingales and asymptotic properties of the free energy of directed polymers in a random environment, Stochastic Processes Appl., Volume 119 (2009) no. 10, pp. 3101-3132 | Zbl | MR
\bibitem{MR2206313} Merlevède, Florence; Peligrad, Magda; Utev, Sergey Recent advances in invariance principles for stationary sequences, Probab. Surv., Volume 3 (2006), pp. 1-36 | MR | Zbl
\bibitem{MR3798239} Peligrad, Magda; Zhang, Na Martingale approximations for random fields, Electron. Commun. Probab., Volume 23 (2018), 28, 9 pages | MR | Zbl
\bibitem{MR2000642} Račkauskas, Alfredas; Suquet, Charles Necessary and sufficient condition for the Lamperti invariance principle, Teor. Ĭmovirn. Mat. Stat., Volume 68 (2003), pp. 115-124 | Zbl
\bibitem{MR2141332} Račkauskas, Alfredas; Suquet, Charles Central limit theorems in Hölder topologies for Banach space valued random fields, Teor. Veroyatn. Primen., Volume 49 (2004) no. 1, pp. 109-125 | DOI | Zbl
\bibitem{MR2054586} Račkauskas, Alfredas; Suquet, Charles Necessary and sufficient condition for the functional central limit theorem in Hölder spaces, J. Theor. Probab., Volume 17 (2004) no. 1, pp. 221-243 | DOI | Zbl
\bibitem{MR2265057} Račkauskas, Alfredas; Suquet, Charles Hölderian invariance principles and some applications for testing epidemic changes, Long memory in economics, Springer, 2007, pp. 109-128 | Zbl | DOI
\bibitem{MR2340883} Račkauskas, Alfredas; Suquet, Charles; Zemlys, Vaidotas A Hölderian functional central limit theorem for a multi-indexed summation process, Stochastic Processes Appl., Volume 117 (2007) no. 8, pp. 1137-1164 | DOI | Zbl
\bibitem{MR606989} Rüschendorf, Ludger Ordering of distributions and rearrangement of functions, Ann. Probab., Volume 9 (1981) no. 2, pp. 276-283 | MR | Zbl
\bibitem{MR79851} Spitzer, Frank A combinatorial lemma and its application to probability theory, Trans. Am. Math. Soc., Volume 82 (1956), pp. 323-339 | DOI | MR | Zbl
\bibitem{MR3427925} Volný, Dalibor A central limit theorem for fields of martingale differences, C. R. Math. Acad. Sci. Paris, Volume 353 (2015) no. 12, pp. 1159-1163 | Numdam | MR | DOI | Zbl
\bibitem{MR3913270} Volný, Dalibor On limit theorems for fields of martingale differences, Stochastic Processes Appl., Volume 129 (2019) no. 3, pp. 841-859 | MR | DOI | Zbl
\bibitem{MR0246359} Wichura, Michael J. Inequalities with applications to the weak convergence of random processes with multi-dimensional time parameters, Ann. Math. Stat., Volume 40 (1969), pp. 681-687 | Zbl | MR | DOI
\bibitem{MR4166203} Zhang, Na; Reding, Lucas; Peligrad, Magda On the quenched central limit theorem for stationary random fields under projective criteria, J. Theor. Probab., Volume 33 (2020) no. 4, pp. 2351-2379 | Zbl | MR | DOI
\end{thebibliography}
\end{document}
